\section{Simulation protocols and statistical conventions}
\label{sec:protocols}

\textit{Simulation.} We use a bit-parallel implementation in which each 64-bit word holds 64
sites and the neighbor counts and flip decisions are computed with bitwise operations. Each site
draws one uniform random number per step and flips if it is below the flip probability of its
class $(\sigma_i, j_i)$; all classes of a word share one bit-serial comparison, so the cost does
not depend on the number of distinct flip probabilities. Random numbers come from the
counter-based Philox$4\times 32$-10 generator~\cite{salmon2011}, keyed by the run's seed and
indexed by (time, row, word), so every run is bit-for-bit reproducible independently of the
number of threads and of restarts from checkpoints. The code passes a bit-exact comparison with a
naive per-site implementation for several rules and lattice shapes, and statistical tests of
the per-class flip frequencies (Supplemental Sec.~\ref{app:numerics}). Near the transition it performs
about $10^9$ site updates per second per core. The extended dense-lattice
campaign uses a CUDA backend built from the same integer update and random-number
functions. CPU and GPU outputs agree bit for bit in the validation cases,
including large lattices, directed profiles, and checkpoints transferred in
both directions (Supplemental Sec.~\ref{app:numerics}).

\textit{Protocols.} We use five kinds of runs.
(i) \emph{Bulk}: random initial condition with density $1/2$ on an $L\times L$ torus.
(ii) \emph{Strip fronts}: on an $L_x\times L_y$ torus, the rows $y<L_y/2$ start random and the
rest dead, giving two flat fronts between the active phase and the absorbing state. The per-row
interface counts are recorded, each front is located at a fixed fraction of
the relaxed interior activity, and the velocity $v$ is half the slope of the
strip width, $W(t) = W_0 + 2vt$,
fitted while the geometry remains a separated pair of fronts. The reference
level and validity windows are specified in Sec.~\ref{app:numerics}.
(iii) \emph{Tilted bands}: a band of direction $(a,b)$, the sites with
$(a y - b x)\bmod L < L/2$, on an $L\times L$ torus. Its two edges have total length
$2L\sqrt{a^2+b^2}$. Profiles binned along the normal locate both edges directly;
their separation gives the normal velocity without converting total activity
into an area.
(iv) \emph{Droplets and seeds}: a disk of radius $R$ whose sites are alive with probability $1/2$
(or a $2\times 2$ block of live sites) in the all-dead state. Many independent replicas are run;
since a site with $j = 0$ never flips, only the bounding box of the active region needs to be
updated, which is exact and makes each replica cheap. A replica stops when it
dies, reaches the lattice edge, or reaches its declared time horizon.
For the droplet experiment, $P_{\rm grow}$ is the probability of reaching
the edge before extinction and before that horizon. In seed-survival curves,
escaped replicas are retained as a separately recorded contribution; they
bound rather than directly measure survival after escape.
(v) \emph{Walls}: all-alive rows $y<L_y/2$ and all-dead rows above, i.e., two flat walls between
the two absorbing phases and no active phase at $t = 0$.
Times are in steps (one synchronous update of all sites). Error bars are standard errors over
independent runs unless stated otherwise.

\textit{Statistical analysis.} The reanalyzed flagship bulk data and the new
fixed-manifest campaigns use every run in complete, predetermined seed groups,
with no density filter or conditioning on survival. Supporting parameter surveys
retain their separately stated windows and selection protocols.
A partially finished fixed group is retained for diagnostics but not
used as an unconditional ensemble: absorption can make a run finish sooner
than a surviving trajectory. We integrate each sampled
trajectory in time before averaging independent runs; a trajectory is extended
by its known constant state after absorption. The half-density time is the first
recorded crossing of $\rho<0.145$ after $t=10^3$, half a reference density
$\rho_{\rm ref}=0.29$ near the flagship transition. It is distinct from complete
consensus. Five legacy $L=64$ runs already below this threshold at the analysis
origin are marked left censored at $10^3$, with lower bound zero. Other
first observed crossings retain the preceding output time as a recording-grid
diagnostic. Activity is not monotone, so brief unobserved dips can occur earlier;
these brackets are not rigorous bounds on its first physical passage.
Unobserved events are right censored, and survival curves and restricted mean times retain
that censoring. An event-free ensemble supplies right-censored observations
and a bound on finite-time survival, not a measured mean lifetime.
Independent-run resampling gives the reported bootstrap
intervals. For the four-run large-system groups these are approximate measures
of observed run variation and do not account for rare outcomes absent from
the sample. Event counts and systematic fitting-window checks are reported
separately.

%=====================================================================================

The percentile intervals for four-run ensembles are displayed as conditional
summaries of observed run variation, not calibrated $95\%$ error guarantees.
We show individual paired changes where their sign is used in the text.
For the two small directional contrasts and midpoint density drift we also
report a small-sample $t$ reference calculation, keeping its normality and
linearization assumptions explicit. Neither method substitutes for more
independent trajectories when precise tail coverage is essential.
